\documentclass[11pt]{article}
\usepackage[margin=1in]{geometry}
\usepackage[pdftex]{graphicx}
\usepackage{amsmath,amssymb,amsthm}
\usepackage{custom}

% plotting
\usepackage{pgfplots}
\pgfplotsset{compat=1.18}

% color boxes
\usepackage[many]{tcolorbox}
\tcbuselibrary{theorems}
\newtcbtheorem{psexample}{Example}{colback=green!10,colframe=green!40!black,fonttitle=\bfseries,breakable,enhanced}{ex}
\newtcbtheorem{psidea}{Idea}{colback=blue!10,colframe=blue!40!black,fonttitle=\bfseries,breakable,enhanced}{id}
\newtcbtheorem{psremark}{Remark}{colback=purple!10,colframe=purple!40!black,fonttitle=\bfseries,breakable,enhanced}{rm}
\newtcbtheorem{pssolution}{Solution}{colback=orange!10,colframe=orange!40!black,fonttitle=\bfseries,breakable,enhanced}{sn}
\tcbset{noparskip/.style={before={\par\pagebreak[0]\medskip\parindent=0pt},
after={\par\medskip}}}

% code for totaling up points
\usepackage{totcount}
\newtotcounter{tpts}
\newtotcounter{cpts}
\newtotcounter{endpts}
\newcommand{\clearpts}{\addtocounter{tpts}{\value{cpts}} \setcounter{cpts}{0}}
\newcommand{\pts}[1]{\clearpts \setcounter{cpts}{#1}}
\newcommand{\totpts}{\setcounter{endpts}{\totvalue{tpts} + \totvalue{cpts}}\arabic{endpts}}

% formatting for problems and solutions
\definecolor{ptred}{rgb}{0.7,0.1,0.1}
\newcommand{\ptfmt}[1]{\textbf{\color{ptred}#1\color{black}}}
\definecolor{advblue}{rgb}{0.1,0.1,0.7}
\newcommand{\advanced}{\!\textbf{\color{advblue}[A]\color{black}}\ }

\newtheorem*{answer}{Answer}

\makeatletter
\newtheoremstyle{mystyle}
{\topsep}               % space above
{\topsep}               % space below
{}                      % bodyfont
{}                      % indent
{\bfseries}             % headfont
{}                      % punctuation
{0.6em}                 % space after head
{\llap{[\ptfmt{\arabic{cpts}}]\hspace{.6em}}\thmname{#1}\thmnumber{ #2}\thmnote{\normalfont{ (#3)}}{\bfseries .}}  %theoremheadspec
\theoremstyle{mystyle}
\newtheorem{pproblem}{Problem}
\makeatother

% multiple-choice lists for the F = ma style questions at the end
\newlist{mcchoices}{enumerate}{1}
\setlist[mcchoices,1]{label=(\Alph*), itemsep=1pt, topsep=3pt, parsep=0pt}

% headers and footers
\usepackage{fancyhdr}
\pagestyle{fancy}
\lhead{Young's Modulus}
\chead{}
\rhead{F = ma Toolkit}
\lfoot{}
\cfoot{\thepage}
\rfoot{}
\renewcommand{\headrulewidth}{0.4pt}
\setlength{\headheight}{14pt}

\newcommand{\psettitle}[1]{
    \begin{center}
    \huge \textbf{#1}
    \end{center}
}

\linespread{1.03} % give a little extra room
\setlength{\parindent}{0.2in} % reduce paragraph indent a bit

% ============================================================
% SOLUTIONS TOGGLE
% Set \showsoltrue to PRINT solutions, \showsolfalse to HIDE them.
% (When building from the command line you can instead pass
%  \def\HIDESOLUTIONS{} to force the student/no-solutions version.)
% ============================================================
\newif\ifshowsol
\showsoltrue                       % <-- flip to \showsolfalse for the student copy
\ifdefined\HIDESOLUTIONS \showsolfalse \fi

\usepackage{environ}

% Solutions are collected as they are met and printed together at the end,
% rather than sitting under their problem: reading a problem should not put its
% answer in your eyeline. Each entry is tagged with the number of the problem it
% belongs to, captured from the pproblem counter at the moment the solution
% appears, so inserting or reordering problems needs no manual renumbering.
\newtoks\soltoks
\soltoks{}
\newcommand{\solentry}[2]{\par\medskip\noindent\textbf{Problem #1.}\quad #2\par}

\ifshowsol
  % The token-register assignment must be global, or the environment's own
  % group discards each entry as soon as the solution ends.
  \NewEnviron{solution}{%
    \begingroup
    \edef\x{\endgroup
      \noexpand\global\noexpand\soltoks{\the\soltoks
        \noexpand\solentry{\thepproblem}{\unexpanded\expandafter{\BODY}}}}%
    \x
  }
  % On a fresh page, so the problem pages are identical in both builds and a
  % student can print only the part they are meant to work from.
  \newcommand{\printsolutions}{\clearpage\section*{Solutions}\the\soltoks}
\else
  \NewEnviron{solution}{}
  \newcommand{\printsolutions}{}
  \NewEnviron{killsol}{}
  \let\answer\killsol
  \let\endanswer\endkillsol
\fi

\begin{document}

\psettitle{Young's Modulus}
\begin{center}
\large How stiff is the \emph{stuff}, as opposed to how stiff is the \emph{thing}
\end{center}

\noindent
Hang a mass from a steel cable and the cable stretches --- a little. Let go and it springs
back. So a cable is a spring, and it must have a spring constant $k$. But $k$ depends on
how long the cable is and how thick it is, which are facts about \emph{that particular
cable}, not about steel. Young's modulus is the number that describes the steel itself.
Everything in this handout follows from separating those two ideas. There is a total of
\ptfmt{\totpts} points; problems marked \advanced are a little harder.

\section{Stress, Strain, and the Definition}

\begin{psidea}{Two ratios, and the number connecting them}{}
Take a rod or wire of length $\ell$ and cross-sectional area $A$, and pull on it with
force $F$ along its length. Define
\[
\text{\textbf{stress}} \quad \sigma = \frac{F}{A},
\qquad\qquad
\text{\textbf{strain}} \quad \varepsilon = \frac{\Delta \ell}{\ell}.
\]
Stress is ``force spread over the area that carries it.'' Strain is ``fractional
stretch'' --- how much longer, as a fraction of what you started with. For small
stretches the two are proportional, and the constant of proportionality is
\textbf{Young's modulus} $Y$:
\[
\boxed{\;\sigma = Y \varepsilon \;}
\qquad\text{that is,}\qquad
\boxed{\;\frac{F}{A} = Y\,\frac{\Delta \ell}{\ell}\;}
\]
This is Hooke's law, rewritten so that only the \emph{material} appears. $Y$ is a property
of steel, or copper, or bone --- not of the particular piece you happen to be holding.
\end{psidea}

\begin{psremark*}{}{}
Why divide by $A$? Because a rope twice as thick shares the load between twice as much
material, so each part of the material is pulled half as hard. Why divide by $\ell$?
Because a rope twice as long has twice as much material to stretch, so the same
\emph{fractional} stretch produces twice the absolute stretch. Stress and strain are
built precisely to cancel out the size of the object.
\end{psremark*}

\pts{3}
\begin{pproblem}
A wire of cross-sectional area $A = 2.0\;\mathrm{mm}^2$ carries a tension of
$F = 400\;\mathrm{N}$.
\begin{subpart}
\item What is the stress, in pascals?
\item If the wire is $1.5\;\mathrm{m}$ long and stretches by $0.60\;\mathrm{mm}$, what is
the strain?
\end{subpart}
\end{pproblem}
\begin{solution}
(a) $A = 2.0\;\mathrm{mm}^2 = 2.0\times10^{-6}\;\mathrm{m}^2$, so
$\sigma = F/A = 400 / (2.0\times 10^{-6}) = 2.0\times 10^{8}\;\mathrm{Pa}
= 200\;\mathrm{MPa}$.

(b) $\varepsilon = \Delta \ell/\ell = (0.60\times10^{-3})/1.5 = 4.0\times 10^{-4}$. Note
this is a pure number --- no units.
\end{solution}

\section{Units, and Why They Are the Units of Pressure}

\begin{psidea}{$Y$ is measured in pascals}{}
Strain is a length divided by a length, so it is \textbf{dimensionless}. That means
$Y = \sigma/\varepsilon$ has exactly the same units as stress:
\[
[Y] = \frac{[F]}{[A]} = \frac{\mathrm{N}}{\mathrm{m}^2} = \mathrm{Pa},
\qquad\text{and in base units}\qquad
[Y] = \frac{MLT^{-2}}{L^2} = M L^{-1} T^{-2}.
\]
Those are the dimensions of pressure, and of energy per unit volume --- all three are the
same thing dimensionally. Real values are large, so $Y$ is usually quoted in gigapascals
($1\;\mathrm{GPa} = 10^9\;\mathrm{Pa}$):
\begin{center}
\begin{tabular}{lr@{\qquad}lr}
rubber & $0.01$--$0.1$ GPa & aluminum & $70$ GPa\\
nylon & $3$ GPa & steel & $200$ GPa\\
wood (along the grain) & $10$ GPa & tungsten & $400$ GPa\\
bone & $15$ GPa & diamond & $1000$ GPa\\
\end{tabular}
\end{center}
\end{psidea}

\pts{2}
\begin{pproblem}
Show, using base dimensions $M$, $L$, $T$, that Young's modulus has the same dimensions as
energy per unit volume. (This is not a coincidence: stretching a wire stores elastic
energy, and $Y$ sets how much.)
\end{pproblem}
\begin{solution}
Energy has dimensions $ML^2T^{-2}$ and volume has dimensions $L^3$, so energy per volume
is $ML^2T^{-2}/L^3 = ML^{-1}T^{-2}$. That is exactly the result computed above for $Y$.
\end{solution}

\begin{psidea}{An intuition for why the number is so big}{}
Put $\varepsilon = 1$ into $\sigma = Y\varepsilon$. You get $\sigma = Y$. So
\emph{Young's modulus is the stress you would need to double the length of the object} ---
if the material kept behaving linearly all the way there, which of course it never does.
Steel snaps at a stress of roughly $0.4\;\mathrm{GPa}$, some $500$ times smaller than its
$Y = 200\;\mathrm{GPa}$. Turned around, that says steel breaks at a strain of about
$1/500 = 0.2\%$. \textbf{Real strains in stiff materials are tiny}, which is why the
modulus is huge.
\end{psidea}

\pts{3}
\begin{pproblem}
Nylon has $Y \approx 3\;\mathrm{GPa}$ and breaks at a stress of about
$75\;\mathrm{MPa}$. At what strain does it break, and how does that compare with steel's
$0.2\%$? Does the more easily broken material stretch more or less before failing?
\end{pproblem}
\begin{solution}
$\varepsilon = \sigma/Y = (75\times10^6)/(3\times10^9) = 0.025$, i.e.\ $2.5\%$ --- about
twelve times the strain steel tolerates. Nylon fails at a much \emph{lower stress} but a
much \emph{higher strain}. ``Strong'' (large breaking stress) and ``stiff'' (large $Y$)
are different properties, and nylon is neither as strong nor as stiff as steel, yet it
stretches far more before it goes.
\end{solution}

\section{From Young's Modulus to a Spring Constant}

\begin{psidea}{$k = YA/\ell$ --- the most useful line in this handout}{}
Rearrange $F/A = Y \Delta\ell/\ell$ to put it in the form $F = k\,\Delta\ell$:
\[
F = \frac{YA}{\ell}\,\Delta \ell
\qquad\Longrightarrow\qquad
\boxed{\;k = \frac{YA}{\ell}\;}
\]
Read this carefully, because it is the whole point:
\begin{itemize}
\item[$\bullet$] $Y$ belongs to the \textbf{material}. Cut a steel cable in half and both
halves still have $Y = 200\;\mathrm{GPa}$.
\item[$\bullet$] $k$ belongs to the \textbf{object}. Cut the cable in half and each half
has \emph{twice} the spring constant of the original.
\item[$\bullet$] \textbf{Longer $\Rightarrow$ floppier.} $k \propto 1/\ell$.
\item[$\bullet$] \textbf{Thicker $\Rightarrow$ stiffer.} $k \propto A$, so $k \propto r^2$
for a circular wire of radius $r$.
\end{itemize}
\end{psidea}

\begin{psremark*}{}{}
You can get both dependences without any algebra, just from how springs combine. A wire
of length $2\ell$ is two wires of length $\ell$ joined end to end --- springs in
\textbf{series}, and $1/k_{\text{tot}} = 1/k + 1/k$ gives $k_{\text{tot}} = k/2$. A wire
of area $2A$ is two wires of area $A$ lying side by side sharing the load --- springs in
\textbf{parallel}, and $k_{\text{tot}} = k + k = 2k$. Both match $k = YA/\ell$ exactly.
\end{psremark*}

\pts{4}
\begin{pproblem}
A steel wire ($Y = 200\;\mathrm{GPa}$) has length $\ell = 2.0\;\mathrm{m}$ and
cross-sectional area $A = 1.0\;\mathrm{mm}^2$. A $20\;\mathrm{kg}$ mass hangs from it.
Take $g = 10\;\mathrm{m/s^2}$.
\begin{subpart}
\item Find the spring constant $k$ of the wire.
\item Find how far it stretches.
\end{subpart}
\end{pproblem}
\begin{solution}
(a) $k = YA/\ell = (200\times10^{9})(1.0\times10^{-6})/2.0 = 1.0\times 10^{5}\;
\mathrm{N/m}$.

(b) The tension is $F = mg = 200\;\mathrm{N}$, so
$\Delta \ell = F/k = 200/(1.0\times10^{5}) = 2.0\times 10^{-3}\;\mathrm{m}
= 2.0\;\mathrm{mm}$.

As a check, the strain is $2.0\;\mathrm{mm}/2.0\;\mathrm{m} = 10^{-3}$, and the stress is
$200\;\mathrm{N}/10^{-6}\;\mathrm{m}^2 = 2\times10^{8}\;\mathrm{Pa}$; their ratio is
$2\times 10^{11}\;\mathrm{Pa} = Y$, as it must be.
\end{solution}

\pts{3}
\begin{pproblem}
Two wires are made of the same metal. Wire 2 is three times as long as wire 1 and has
twice the radius. By what factor is wire 2's spring constant larger or smaller than wire
1's? By what factor does their Young's modulus differ?
\end{pproblem}
\begin{solution}
$k = YA/\ell \propto r^2/\ell$. Doubling $r$ multiplies $k$ by $4$; tripling $\ell$
divides it by $3$. So $k_2/k_1 = 4/3$. Young's modulus is identical for both --- they are
the same metal --- so it differs by a factor of $1$.
\end{solution}

\pts{3}
\begin{pproblem}
A rubber band and a steel wire of identical dimensions are each stretched by the same
\emph{force}. Which stretches further, and by roughly what factor? (Use
$Y_{\text{rubber}} \approx 0.05\;\mathrm{GPa}$.)
\end{pproblem}
\begin{solution}
Same $F$, same $A$ means same stress, so $\varepsilon = \sigma/Y$ says the strain is
inversely proportional to $Y$. The rubber stretches by a factor
$Y_{\text{steel}}/Y_{\text{rubber}} = 200/0.05 = 4000$ more than the steel.
\end{solution}

\section{Oscillations on a Wire}

\begin{psidea}{A hanging mass on a wire is just a mass on a spring}{}
Once you know $k = YA/\ell$, vertical oscillations need nothing new:
\[
f = \frac{1}{2\pi}\sqrt{\frac{k}{m}}
= \frac{1}{2\pi}\sqrt{\frac{YA}{\ell m}}.
\]
Gravity does not appear. It shifts the equilibrium point down by $mg/k$ but does not
change the stiffness, exactly as for an ordinary vertical spring.
\end{psidea}

\begin{psexample}{The question you saw on the practice exam}{}
\emph{Steel has $Y = 200\;\mathrm{GPa}$. A mass $m = 300\;\mathrm{kg}$ hangs from a steel
wire of circular cross-section with radius $r = 0.15\;\mathrm{cm}$ and length
$\ell = 2\;\mathrm{m}$. Find the frequency of vertical oscillations.}

First the area: $r = 0.15\;\mathrm{cm} = 1.5\times 10^{-3}\;\mathrm{m}$, so
\[
A = \pi r^2 = \pi (1.5\times10^{-3})^2 = 7.07\times 10^{-6}\;\mathrm{m}^2.
\]
Then the spring constant:
\[
k = \frac{YA}{\ell} = \frac{(200\times10^{9})(7.07\times10^{-6})}{2}
= 7.07\times 10^{5}\;\mathrm{N/m}.
\]
Then the frequency:
\[
f = \frac{1}{2\pi}\sqrt{\frac{k}{m}}
= \frac{1}{2\pi}\sqrt{\frac{7.07\times10^{5}}{300}}
= \frac{1}{2\pi}\sqrt{2360}
= \frac{48.6}{2\pi} \approx 7.7\;\mathrm{Hz}.
\]
The only step that is specific to this topic is $k = YA/\ell$. Everything after it is
ordinary SHM. Watch the centimetres: writing $r = 0.15$ instead of $0.0015$ costs a
factor of $100$ in $r$, hence $10^4$ in $A$ and $100$ in $f$.
\end{psexample}

\pts{3}
\begin{pproblem}
In the example above, how far below the wire's unstretched end does the mass hang at
equilibrium? Compare that sag with the amplitude of a small oscillation, and comment on
whether the wire ever goes slack.
\end{pproblem}
\begin{solution}
The static stretch is $\Delta \ell = mg/k = (300)(9.8)/(7.07\times 10^5)
= 4.2\times 10^{-3}\;\mathrm{m} \approx 4.2\;\mathrm{mm}$. Any oscillation with amplitude
smaller than $4.2\;\mathrm{mm}$ keeps the wire in tension the whole time, so the motion
stays simple harmonic. A larger amplitude would let the wire go slack at the top of the
motion, where it can push nothing --- and then the motion is no longer SHM.
\end{solution}

\section{Dimensional Analysis with $Y$}

\begin{psidea}{$\sqrt{Y/\rho}$ is a speed}{}
$Y$ has dimensions $ML^{-1}T^{-2}$ and density $\rho$ has dimensions $ML^{-3}$. Divide:
\[
\left[\frac{Y}{\rho}\right] = \frac{ML^{-1}T^{-2}}{ML^{-3}} = L^2T^{-2},
\]
which is a velocity squared. So $\sqrt{Y/\rho}$ is the only speed you can build from a
material's stiffness and its density --- and it is, up to a factor of order $1$, the
\textbf{speed of sound in that material}. Stiffer material, faster sound; heavier
material, slower sound.
\end{psidea}

\pts{3}
\begin{pproblem}
Estimate the speed of sound in steel, using $Y = 200\;\mathrm{GPa}$ and
$\rho = 8000\;\mathrm{kg/m^3}$. Compare with the speed of sound in air
($340\;\mathrm{m/s}$).
\end{pproblem}
\begin{solution}
$v = \sqrt{Y/\rho} = \sqrt{(2\times10^{11})/(8\times10^{3})} = \sqrt{2.5\times 10^{7}}
= 5.0\times 10^{3}\;\mathrm{m/s}$. That is about $15$ times the speed of sound in air.
(The measured value for steel is about $5900\;\mathrm{m/s}$, so the estimate is good.)
\end{solution}

\pts{4}
\begin{pproblem} \advanced
A wire of length $\ell$, area $A$, density $\rho$ and Young's modulus $Y$ hangs with a
mass $m$ on the end. Without solving anything, use dimensional analysis to argue that the
oscillation frequency must have the form
$f = \dfrac{1}{2\pi}\sqrt{\dfrac{YA}{\ell m}} \times g(\text{dimensionless})$, and
identify a dimensionless ratio the leftover function could depend on.
\end{pproblem}
\begin{solution}
$YA$ has dimensions $(ML^{-1}T^{-2})(L^2) = MLT^{-2}$, a force. Dividing by $\ell m$ gives
$MLT^{-2}/(LM) = T^{-2}$, so $\sqrt{YA/(\ell m)}$ is a frequency --- the only one you can
build from those four quantities. Anything else must be dimensionless. The natural
dimensionless ratio available is the mass of the wire to the hanging mass,
$\rho A \ell / m$; the exact answer does depend on it slightly, because the wire itself
has to move too. When the wire is light compared with the load, that ratio is near zero
and $g \to 1$, recovering the formula used above.
\end{solution}

\section{Common Traps}

\begin{psremark*}{}{}
\begin{itemize}
\item[$\bullet$] \textbf{``Longer means less stiff'' applies to $k$, never to $Y$.} A
longer rope has a smaller spring constant and the identical Young's modulus. If a problem
changes the length and you find yourself changing $Y$, stop.
\item[$\bullet$] \textbf{Watch the area units.} $1\;\mathrm{mm}^2 = 10^{-6}\;\mathrm{m}^2$
and $1\;\mathrm{cm}^2 = 10^{-4}\;\mathrm{m}^2$. Squaring turns a small slip into a factor
of $100$ or $10^4$.
\item[$\bullet$] \textbf{A radius given in centimetres must be converted before squaring.}
This is the single most common way to lose the steel-wire problem.
\item[$\bullet$] \textbf{Strain is a fraction, not a length.} $\Delta \ell$ is a length;
$\Delta \ell / \ell$ is the strain. Only the second one goes into $\sigma = Y\varepsilon$.
\item[$\bullet$] \textbf{Stiff is not the same as strong.} $Y$ tells you how much it
stretches under load; the breaking stress tells you when it snaps. Materials can be high
in one and low in the other.
\end{itemize}
\end{psremark*}

\section{F = ma Style Questions}

\noindent
Four questions in the style of the exam. Answer each in about two minutes.

\pts{4}
\begin{pproblem}
A steel wire of length $\ell$ and cross-sectional area $A$ stretches by $1.0\;\mathrm{mm}$
when a certain mass is hung from it. A second wire, made of the same steel with the same
cross-sectional area but \emph{twice the length}, carries the same mass. Compared with the
first wire, the second wire has
\begin{mcchoices}
\item the same Young's modulus, half the spring constant, and stretches $2.0\;\mathrm{mm}$.
\item half the Young's modulus, half the spring constant, and stretches $2.0\;\mathrm{mm}$.
\item the same Young's modulus, twice the spring constant, and stretches
$0.50\;\mathrm{mm}$.
\item the same Young's modulus, the same spring constant, and stretches
$1.0\;\mathrm{mm}$.
\item twice the Young's modulus, the same spring constant, and stretches
$0.50\;\mathrm{mm}$.
\end{mcchoices}
\end{pproblem}
\begin{solution}
\textbf{(A).} Young's modulus is a property of steel and does not know how long the wire
is, so it is unchanged. The spring constant $k = YA/\ell$ is halved when $\ell$ doubles.
With the same force and half the spring constant, $\Delta \ell = F/k$ doubles to
$2.0\;\mathrm{mm}$.

Equivalently and even faster: the same mass on the same cross-section gives the same
stress, hence the same \emph{strain}. Same fractional stretch on twice the length is twice
the absolute stretch.
\end{solution}

\pts{4}
\begin{pproblem}
A steel wire ($Y = 200\;\mathrm{GPa}$) of length $2.0\;\mathrm{m}$ and cross-sectional
area $1.0\;\mathrm{mm}^2$ supports a $20\;\mathrm{kg}$ mass. Taking
$g = 10\;\mathrm{m/s^2}$, the wire stretches by about
\begin{mcchoices}
\item $0.20\;\mathrm{mm}$
\item $0.50\;\mathrm{mm}$
\item $1.0\;\mathrm{mm}$
\item $2.0\;\mathrm{mm}$
\item $4.0\;\mathrm{mm}$
\end{mcchoices}
\end{pproblem}
\begin{solution}
\textbf{(D).} The tension is $F = mg = 200\;\mathrm{N}$, so the stress is
$\sigma = 200/(1.0\times10^{-6}) = 2.0\times 10^{8}\;\mathrm{Pa}$. The strain is
$\varepsilon = \sigma/Y = (2.0\times10^{8})/(2.0\times10^{11}) = 1.0\times 10^{-3}$.
Therefore $\Delta \ell = \varepsilon \ell = (1.0\times10^{-3})(2.0\;\mathrm{m})
= 2.0\times 10^{-3}\;\mathrm{m} = 2.0\;\mathrm{mm}$.
\end{solution}

\pts{4}
\begin{pproblem}
The speed $v$ of a sound pulse traveling along a thin solid rod depends only on the rod's
Young's modulus $Y$ and its density $\rho$. For steel, $Y = 200\;\mathrm{GPa}$ and
$\rho = 8000\;\mathrm{kg/m^3}$. The speed is closest to
\begin{mcchoices}
\item $160\;\mathrm{m/s}$
\item $1600\;\mathrm{m/s}$
\item $5000\;\mathrm{m/s}$
\item $25000\;\mathrm{m/s}$
\item $2.5\times 10^{7}\;\mathrm{m/s}$
\end{mcchoices}
\end{pproblem}
\begin{solution}
\textbf{(C).} Only one combination of $Y$ and $\rho$ has dimensions of speed. Since
$[Y] = ML^{-1}T^{-2}$ and $[\rho] = ML^{-3}$, the ratio $Y/\rho$ has dimensions
$L^2T^{-2}$, so $v = \sqrt{Y/\rho}$. Numerically
$v = \sqrt{(2\times10^{11})/(8\times10^{3})} = \sqrt{2.5\times10^{7}}
= 5.0\times 10^{3}\;\mathrm{m/s}$.

Distractor (E) is $Y/\rho$ without the square root, and (A) is $\sqrt{\rho/Y}$ scaled ---
both are what you get by taking the combination in the wrong order, which the dimensions
rule out immediately.
\end{solution}

\pts{4}
\begin{pproblem}
A mass $m$ hangs from a steel wire of length $\ell$ and cross-sectional area $A$, and
oscillates vertically with frequency $f$. The wire is replaced by one of the same steel
and the same length, but with \emph{four times} the cross-sectional area. The new
frequency is
\begin{mcchoices}
\item $f/4$
\item $f/2$
\item $f$
\item $2f$
\item $4f$
\end{mcchoices}
\end{pproblem}
\begin{solution}
\textbf{(D).} The spring constant is $k = YA/\ell$, so quadrupling $A$ quadruples $k$.
Since $f = \frac{1}{2\pi}\sqrt{k/m}$ and $m$ is unchanged, $f \propto \sqrt{k}$, and
$\sqrt{4} = 2$. The frequency doubles.

Note that $Y$ never changes --- it is the same steel throughout. Only the geometry moved.
\end{solution}

\vspace{1em}
\begin{center}
\emph{End of handout. If you remember only one line, make it $k = YA/\ell$: the material
contributes $Y$, the object contributes $A$ and $\ell$, and everything else is the spring
physics you already know.}
\end{center}

\printsolutions

\end{document}
